% s13 The main theorem for King's printed state
\section{The main theorem for King's printed state}

We state the main theorem in full and prove it by assembling Sections~4--12. The object is King's printed state $\ZKinf(f,0)=\lim_J\lim_k\avK{e^{F(f)}}$ on King's RG cube (\cref{def:s02:Kcube}), block factor $5$, at a fixed photon mass $\lmu>0$ and all positive couplings $e,\lambda,m$ (\cref{lem:s14:CW}). Both sectors are proved on King's convention $\NDK$, on which the existence theorems \cref{hyp:s03:EC,hyp:s03:W1} are posed. In the Wick sector (Leg~W), proved in its tube form, Form~(B), the tube statements \cref{thm:s11:W,lem:s11:TWtube,thm:s08:OS2WT} are written on the comparison convention $\NDa$ (\cref{def:s02:acube}) as templates; the two of them that the interlacing of \cref{lem:s04:G} does not transfer, the bound \eqref{eq:s11:W} and the fixed-cube limits on the tube for the K-family functionals $Z^{a,K}_{k,J}$, are proved directly on King's convention in \cref{thm:s13:TWK}, after which the Wick sector follows on $\NDK$ (\cref{prop:s13:WK}); the comparison family enters through its LIM states, which coincide with King's tube object (\cref{prop:s13:LIMK}). The torus inputs are \cref{hyp:s03:T1,hyp:s03:T2}, and the coupling condition is \Tthree\ (\cref{hyp:s03:T3}), the coupling region of which contains the explicit regions of \cref{prop:s14:T3}; $\cK=\{(\mathrm if,g):g\ge0\}$ is the cone. \Cref{thm:s01:main} is the theorem below.

\begin{theorem}[the axioms for King's printed state at fixed photon mass]\label{thm:s13:main}
Let $e>0$, $\lambda,m,\lmu>0$, $\eps_k=5^{-k}$, $\CJ=[-MJ,MJ)^3$ with King's block integer $M\ge1$, and let $\NDK_{k,J}$ be King's cube state.

\smallskip\noindent\textbf{A. The object.} For every finite family $f=(f_1,\dots,f_N)$ of real $2$-forms in $C^\infty_c(\R^3;\Lambda^2)$ the limit
\[
\ZKinf\Big(\sum_jz_jf_j,0\Big)=\lim_{J\to\infty}\lim_{k\to\infty}\avK{e^{\sum_jz_jF_{\Lk}(f_j)}}
\]
exists locally uniformly in $z\in\C^N$, is entire, depends on the family only through the total source, is real and $\ge1$ at real $z$, lies in $[0,1]$ at purely imaginary $z$, and satisfies $|\ZKinf(\sum_jz_jf_j,0)|\le e^{\norm{g(z)}_2^2}$, $g(z)=\sum_j(\Re z_j)f_j$, with no dependence on $\lmu$, $e$, $\lambda$, $m$ or $M$.

\smallskip\noindent\textbf{B. Leg~C (the curvature field).} The Schwinger functions $S^K_n(f_1,\dots,f_n)=\partial_{z_1}\cdots\partial_{z_n}\ZKinf(\sum_jz_jf_j,0)|_{z=0}$ satisfy:
\begin{enumerate}[label=\textup{(\roman*)},leftmargin=*]
\item \emph{\OSz.} $|\ZKinf(\sum_jz_jf_j,0)|\le e^{\norm{g(z)}_2^2}$ for all $z\in\C^N$, and also $\le\exp(\frac12\lmu^{-2}\norm{\sum_j(\Re z_j)\dstar f_j}_2^2)$; the same two bounds, with the lattice norms $\norm{\cdot}_{2,k}$ and $\dstar_k$, hold at every $(k,J)$ for the lattice functionals $\avK{e^{\sum_jz_jF_{\Lk}(f_j)}}$. For real test functions
\[
|S^K_n(f_1,\dots,f_n)|\le(\sqrt2\,e)^n(n!)^{1/2}\prod_{i=1}^n\norm{f_i}_2\le(e\pi/2)^n(n!)^{1/2}\prod_{i=1}^n\fournorm{f_i};
\]
each $S^K_n$ extends to a continuous symmetric multilinear form on $\cS(\R^3;\Lambda^2_\C)^n$, a tempered distribution, the bound carrying the factor $2^n$ for complex test functions; so \Ezz\ holds with $s=4$, $\sigma_n=(e\pi)^n(n!)^{1/2}$ and every $\beta>\frac12$.
\item \emph{\OSo.} The $S^K_n$ are invariant under all translations of $\R^3$ and under the cubic group $\Oh$ about every point, $F$ transforming as a $2$-form. Moreover $\ZKinf$ is invariant under the Euclidean group $\Euc$ and the $S^K_n$ are Euclidean covariant.
\item \emph{\OSt.} Reflection positivity about every coordinate plane $\{x_\nu=c\}$, $\nu\in\{1,2,3\}$, $c\in\R$: for finite families $\zeta_1,\dots,\zeta_N$ of real $2$-form sources supported in $\{x_\nu>c\}$ with arbitrary complex coefficients, $\sum_{i,j}\bar\gamma_i\gamma_j\ZKinf(\theta\zeta_i^*\sqcup\zeta_j,0)\ge0$ for all $\gamma\in\C^N$ (the matrix form for exponentials), and $S(\theta(P)P)\ge0$ for every polynomial $P$ in the smeared curvature with sources supported in $\{x_\nu>c\}$ (the polynomial form). The same holds about every plane $\{x\cdot n=c\}$, $n$ a unit vector.
\item \emph{\OSth.} $S^K_n$ is real for real test functions, symmetric, and zero for odd $n$.
\item \emph{The measure.} For every even $r\ge14$ there is exactly one Borel probability measure $\mu^{K,H}$ on $\cH^{-r}(\R^3;\Lambda^2)$ whose moments on $C^\infty_c$ are the $S^K_n$; its characteristic functional is $\ZKinf(\mathrm i\,\cdot,0)$, extended by continuity to $\cS$; its image $\mu^K_S$ on $H^{-r}_{\rm loc}(\R^3;\Lambda^2)$ is the unique Borel probability measure there with this characteristic functional, and the unique one with these moments. For every real $\chi\in C^\infty_c(\R^3)$ and every box $Q$ containing $\supp\chi$ in its interior, the laws of the localised lattice curvature fields $\chi\Phi^K_{k,J}$ under $\NDK_{k,J}$ on $H^{-r}_{\rm per}(Q;\Lambda^2)$ converge weakly, along the full sequence $k\to\infty$ at every fixed $J\ge J_0(\chi)$ and then as $J\to\infty$, to the law of $\chi\Phi$ under $\mu^K_S$. Nothing is asserted about the laws of the unlocalised lattice fields at fixed $J$.
\item \emph{Weak \OSf.} $\mu^K_S$ is $\Cub$-invariant and is an integral $\int\mu_\xi\,d\nu(\xi)$ of translation-ergodic, reflection-positive, $\theta_j$-invariant components which cluster in the mean and under diagonal separation, have a unique joint vacuum in each coordinate direction, and satisfy the factorial growth bound $|S^\xi_N(u_1,\dots,u_N)|\le\sigma^\xi_N\prod_i\norm{u_i}_{r+2}$, $\sigma^\xi_N\le\alpha_\beta(\xi)(N!)^\beta$ for every $\beta>1$, with component constants (no tube bound); there is an $\Oh$-invariant decomposition into mixtures of $k\in\{1,2,3,6\}$ ergodic measures; and $\nu$-almost every component satisfies \OSf\ in the weak form (ergodicity under $\R^3$). Moreover $\nu$-almost every component is $\Euc$-invariant, reflection positive about every plane, and satisfies \OSf\ in each of its forms for every time direction, with $\ker H_n=\C\Omega_n$ in every direction $n$; its moments are real and symmetric. The vanishing of a component's odd moments is not claimed.
\end{enumerate}

\smallskip\noindent\textbf{C. Leg~W, Form~(B), on King's convention.} Let the couplings lie in the coupling region of \cref{hyp:s03:T3}, the set on which \Tthree\ holds, which contains the explicit regions of \cref{prop:s14:T3}, and fix an admissible scheme $a$ (\cref{def:s03:adm}; $z=a^\sharp$ the reference scheme). Let $\sigma_+=s_1$, $H=s_1/16$, $\kappa=C_U/s_1$, $K_-$ be the constants of \eqref{eq:s11:constants} ($K_-$ built from $B_{\rm cell}$; for $a\ne z$ with $C_U$ replaced by $C_U+s_1\sup_k|z_k-a_k|$), $\sigma_0=H/4$ and $K'=\kappa\pi^2+K_-$. Then for every finite family $\cF=(f_1,\dots,f_n;h_1,\dots,h_p)$ of real $2$-forms and real scalars the tube object
\[
G^K_\infty(\zeta)=\lim_{J\to\infty}\lim_{k\to\infty}Z^{a,K}_{k,J}\Big(\sum_iz_if_i,\sum_jw_jh_j\Big),\qquad\zeta=(z,w)\in\Dtube=\{\norm{k_+(\zeta)}_\infty<\sigma_+/2,\ \norm{k_-(\zeta)}_\infty<H/2\},
\]
exists locally uniformly on the tube and is holomorphic there; on the Wick tube of \cref{def:s08:enlarged} the limit $Z^{a,K}_{\infty,T}=\lim_J\lim_kZ^{a,K}_{k,J}$ exists and equals $G^K_\infty$ where both are defined. The joint Schwinger functions $S_{n,p}(f_1,\dots,f_n;h_1,\dots,h_p)=\partial^{(n)}_z\partial^{(p)}_wG^K_\infty(0)$ of the pair $(F,W^a)$ satisfy:
\begin{enumerate}[label=\textup{(\roman*)},leftmargin=*]
\item \emph{\OSz.} $|G^K_\infty(\zeta)|\le\exp(\norm{g(\zeta)}_2^2+\kappa\norm{k_+(\zeta)}_{L^1}+K_-\wnorm{k_-(\zeta)})$ on $\Dtube$, and the same bound with lattice norms holds for $Z^{a,K}_{k,J}$ uniformly in $k\ge0$ and $J\ge J_1(R)$ when the supports lie in $Q_R$; for real test functions
\[
|S_{n,p}(f_1,\dots,f_n;h_1,\dots,h_p)|\le e^{K'\sigma_0}(\sqrt2\,e)^n(n!)^{1/2}(e/\sigma_0)^p\,p!\prod_{i=1}^n\norm{f_i}_2\prod_{j=1}^p\wnorm{h_j},
\]
with the factor $2^{n+p}$ for complex test functions; the $S_{n,p}$ are tempered distributions and \Ezz\ holds with $s=4$ and every $\beta>1$. The constants $\sigma_+,H,\kappa,K_-$ depend on $(e,\lambda,m,\lmu,M,a)$, on the fixed auxiliary choices $M_1$, $\chi_0$ and the amplitude $s_1\in(0,\sigma_1)$ of \Itwo, and on nothing else; they are not uniform in $\lmu$.
\item \emph{\OSo.} $G^K_\infty(\tau_bR\cdot\zeta)=G^K_\infty(\zeta)$ on $\Dtube$ for every $b\in\R^3$ and $R\in\Oh$, $F$ transforming as a $2$-form and $W^a$ as a scalar (the transformed family has the same tube); hence the $S_{n,p}$ are invariant under translations and under $\Oh$ about every point. No Euclidean clause is asserted.
\item \emph{\OSt.} For finite families $\zeta_1,\dots,\zeta_N$ of real sources supported in $\{x_\nu>c\}$ with complex coefficients of amplitude $\amp(\zeta_i)<\sigma_W\defeq\sigma_+/8$ (in particular all purely imaginary Wick coefficients with arbitrary complex curvature coefficients), $\sum_{i,j}\bar\gamma_i\gamma_jZ^{a,K}_{\infty,T}(\theta\zeta_i^*\sqcup\zeta_j)\ge0$ about every coordinate plane; and the polynomial form $S(\theta(P)P)\ge0$ for polynomials in the smeared $F$ and $W^a$ with sources in $\{x_\nu>c\}$.
\item \emph{\OSth.} The $S_{n,p}$ are symmetric multilinear in the $f$'s and in the $h$'s, real for real test functions, and zero for odd $n$.
\item[\textup{(v),(vi)}] \emph{The measure and weak \OSf.} Clauses (v) and (vi) of Leg~C hold for the joint field on $\cH^{-r}(\R^3;\Lambda^2\oplus\R)$ and $H^{-r}_{\rm loc}(\R^3;\Lambda^2\oplus\R)$, every even $r\ge14$, with moments $S_{n,p}$, characteristic functional $(u,v)\mapsto G^K_\infty(\mathrm iu,\mathrm iv)$ on real test functions (purely imaginary coefficients have $k(\zeta)=0$ and lie in the tube), the component growth bound with the constants of (i), and the local convergence of the laws of the scheme-shifted lattice fields $\chi(F_{\Lk},W^a_{\Lk})$ under $\NDK_{k,J}$ along the full sequence $k\to\infty$ at every fixed $J\ge\max(J_0(\chi),J_1(R_\chi))$, then $J\to\infty$. No Euclidean clause is asserted. Two admissible schemes $a,a'$ give objects differing by the shift $W\mapsto W-\kappa_{a'a}\int(\cdot)$ of the Wick component, $\kappa_{a'a}=\lim_k(a'_k-a_k)$; joint cumulants of order $\ge2$ are scheme-free.
\end{enumerate}
No bound at positive scalar sources beyond the cone $\sup q\le s_1$ of \cref{cor:s11:W2prime} is used.

\smallskip\noindent\textbf{D. The comparison family.} Under the conditions of~C, every increasing sequence of spacings has a continuum subsequence $\cK$ for the Wick sector of the comparison family $\NDa$ (\cref{def:s12:contsub}), and for every such $\cK$ and every family the LIM state $Z^{a,\cK}_\infty=\lim_J\lim_{k\in\cK}Z^a_{k,J}$ exists on $\Dtube$ and equals $G^K_\infty$ (\cref{prop:s13:LIMK}). Hence clauses (i)--(vi) of~C hold for the Wick-sector LIM states of the comparison family, which coincide with King's tube object and do not depend on $\cK$; in (v) the laws of the localised lattice fields $\chi(F_{\Lk},W^a_{\Lk})$ under $\NDa_{k,J}$ converge along $\cK$ at every fixed $J$ with $MJ>R(\chi)+\ell_c$, then as $J\to\infty$, to the law of the joint field under the measure of~C(v).
\end{theorem}

\begin{proof}
\emph{A.} \Cref{thm:s05:CK}(b): by \cref{hyp:s03:EC} the two limits exist locally uniformly on $\C^N$ and are entire, $|\ZKJ|,|\ZKinf|\le e^{\norm{g(z)}_2^2}$, $J\mapsto\ZKJ(\mathrm it)$ is nondecreasing with values in $[0,1]$, and $\ZKinf(x)\ge1$ at real $x$; realness at real $z$ is the symmetry $(A,\vartheta)\mapsto(-A,-\vartheta)$ of every $\NDK_{k,J}$, closed under limits. Each lattice functional depends on the family only through its total source, hence so do the limits. The bound is absolute (\cref{rem:s05:absolute}).

\emph{B.} (i) is \cref{thm:s05:CK}: the finite-volume bound (a) with $\min\{\norm{g}^2_{2,k},\frac12\lmu^{-2}\norm{\dstar_kg}^2_{2,k}\}$ at every $(k,J)$; the two limit bounds (b), the second being (b)'s second bound, carried from the lattice by $\norm{\dstar_kg}_{2,k}\to\norm{\dstar g}_2$, \eqref{eq:s02:riemann}; the Schwinger bounds, the extension to $\cS$, the factor $2^n$ and \Ezz\ (c). (ii): translations are \cref{cor:s06:TK} (the lattice sandwich \cref{lem:s06:TK}, the two limits, the identity theorem, density); the cubic group about every point is \cref{thm:s06:S1K}; the Euclidean clause is \cref{thm:s06:O3}, that is \cref{rot:thm:curv} (the rotation theorem \cref{hyp:s03:RC}) with \cref{cor:s06:TK}. (iii): the matrix form about every coordinate plane and the polynomial form are \cref{lem:s08:OS2K} with \cref{prop:s08:poly}; a plane $\{x\cdot n=c\}$ is the image of a coordinate plane under a Euclidean motion $\gamma$, and by (ii) $\ZKinf(\theta_\gamma\zeta_i^*\sqcup\zeta_j,0)=\ZKinf(\gamma^{-1}\cdot(\theta_\gamma\zeta_i^*\sqcup\zeta_j),0)$ is a Gram entry about the coordinate plane for the sources $\gamma^{-1}\cdot\zeta_i$, supported in the corresponding coordinate half-space. (iv) is \cref{prop:s08:OS3} for $\ZKinf(\cdot,0)$ (through \cref{lem:s08:OS2K}; charge conjugation gives the vanishing for odd $n$). (v) is \cref{thm:s09:S2K}(a), including the local convergence of the lattice laws along the full sequence, with \cref{rem:s09:S2K} for the unlocalised fields; (vi) is \cref{thm:s09:S2K}(b) and (c), whose $O(3)$ input is (ii)'s Euclidean clause; the disclaimer on odd moments is that of \cref{thm:s09:S2K}(c).

\emph{C.} Under \Tthree, \cref{prop:s10:U} gives \Itwo\ with amplitude $s_1\in(0,\sigma_1)$ and constant $C_U$ from the torus inputs \cref{hyp:s03:T1,hyp:s03:T2}; for $a\ne z$ the difference $z_k-a_k$ converges by admissibility, so $\sup_k|z_k-a_k|<\infty$ and $C_U$ is replaced as stated (\cref{hyp:s12:SCH}); \Ithree\ is \cref{prop:s10:RP}, the correlation inequalities are \cref{lem:s02:BFS}, and the cell constant $B^{(a)}_{\rm cell}$ is finite (\cref{hyp:s12:Bcell}). \emph{(i)}: \cref{thm:s13:TWK} gives the bound \eqref{eq:s13:WK} for every $k\ge0$ and $J\ge J_1(R)$, the fixed-cube limits $G^K_J$ and the cube limit $G^K_\infty$ locally uniformly on $\Dtube$, their holomorphy, the bound in the limit, the Schwinger bounds, the factor $2^{n+p}$ and \Ezz, with the stated dependence of the constants and no uniformity in $\lmu$. On the Wick tube $T_{\cF'}$ of the enlarged family of \cref{def:s08:enlarged}, whose scalar members are nonnegative and whose points have $2\norm{k_1}_\infty<\sigma_+/2$, the fixed-cube limits of $Z^{a,K}_{k,J}$ exist locally uniformly by \cref{thm:k02:W1}(i), whose domain $\{\sup_x2k_1<s_W\}$, $s_W=s_1$, contains $T_{\cF'}$, for the reference scheme and hence for $a$ (\cref{def:s03:adm}); their cube limit $Z^{a,K}_{\infty,T}$ is \cref{prop:s08:T}(a) run on King's convention (\cref{prop:s13:WK}(c)), and it equals $G^K_\infty$ on $T_{\cF'}\cap\Dtube$ as the same $J$-limit of the same fixed-cube limits. \emph{(ii)}: \cref{prop:s13:WK}(b). \emph{(iii)}: \cref{prop:s13:WK}(c), that is \cref{thm:s08:OS2WT} with \cref{prop:s08:T} run on King's convention, and the polynomial form \cref{prop:s08:poly}. \emph{(iv)}: \cref{prop:s08:OS3} for $Z^{a,K}_{\infty,T}$ and $G^K_\infty$ (\cref{prop:s13:WK}(c)). \emph{(v),(vi)}: \cref{prop:s13:WK}(d). The Gram points of (v)'s characteristic functional have purely imaginary coefficients and lie in the tube, and the reflection positivity consumed there is (iii) at amplitude $0$ (\cref{rem:s13:amplitudes}). \emph{The scheme clause}: for admissible $a,a'$, $Z^{a',K}_{k,J}(\zeta)=e^{-(a'_k-a_k)\sum_jw_j\sS_{\Lk}(h_j)}Z^{a,K}_{k,J}(\zeta)$ (\cref{def:s02:sources}), where $a'_k-a_k=(z_k-a_k)-(z_k-a'_k)\to\kappa_{a'a}$ and $\sS_{\Lk}(h_j)\to\int h_j$; the factor converges locally uniformly in $\zeta$, so the fixed-cube limits and then the cube limits satisfy $G^{K,a'}_\infty(\zeta)=e^{-\kappa_{a'a}\sum_jw_j\int h_j}G^{K,a}_\infty(\zeta)$ on the tube, the shift $W\mapsto W-\kappa_{a'a}\int(\cdot)$ of the Wick component; a factor $e^{\ell(w)}$ with $\ell$ linear changes only the first cumulants in $w$, so the joint cumulants of order $\ge2$ are scheme-free.

\emph{D.} \Cref{prop:s13:LIMK}; the clauses of~C then hold for $Z^{a,\cK}_\infty=G^K_\infty$, and the convergence of the laws of the localised comparison-cube fields along $\cK$ is clause (v) of the Wick sector of \cref{thm:s12:LIM}, whose hypotheses \Bcell\ and \SCH\ hold for every admissible scheme under \Tthree\ (\cref{hyp:s12:Bcell,hyp:s12:SCH}).
\end{proof}

\begin{remark}[the Euclidean clause]\label{rem:s13:rotation}
The Euclidean clauses of B(ii), (iii), (vi) rest on the rotation theorem \cref{hyp:s03:RC}, proved in the rotation part (\cref{rot:thm:main}): the walled tori on the two lattices have the same continuum limits (\cref{rot:hyp:K}), King's printed state is the lower envelope of the walled tori (\cref{rot:hyp:TIGHT}), and that closing estimate comes from the heavy-shell estimates on King's block-spin flow (\cref{rot:hyp:HS}); the Wick sector carries no Euclidean clause (\cref{thm:s23:FormA}).
\end{remark}

\begin{theorem}[the Wick tube on King's convention]\label{thm:s13:TWK}
Let the couplings satisfy \Tthree\ and let $a$ be an admissible scheme. Let $\sigma_+=s_1$, $H$, $\kappa$ be the constants of \eqref{eq:s11:constants} and let $K_-$ be the constant there, computed from the cell constant $B^{(a)}_{\rm cell}$ of the comparison convention (\cref{hyp:s12:Bcell}), which is finite. Then for every finite family of real sources with supports in $Q_R$, every $k\ge0$, every $J\ge J_1(R)$ and every $\zeta\in\Dtube$,
\begin{equation}\label{eq:s13:WK}
|Z^{a,K}_{k,J}(\zeta)|\ \le\ Z^{a,K}_{k,J}(\Re\zeta)\ \le\ \exp\big(\norm{g(\zeta)}_{2,k}^2+\kappa\norm{k_+(\zeta)}_{1,k}+K_-\wnorm{k_-(\zeta)}\big),
\end{equation}
the bound \eqref{eq:s11:W} for the K-family functionals with the constants of the comparison convention; and the fixed-cube limits $G^K_J(\zeta)=\lim_{k\to\infty}Z^{a,K}_{k,J}(\zeta)$ exist for every $\zeta\in\Dtube$, locally uniformly, so that \cref{thm:s11:W}(i),(ii) hold on $\NDK$: $G^K_J$ and $G^K_\infty=\lim_JG^K_J$ exist locally uniformly on $\Dtube$ and are holomorphic, the bound \eqref{eq:s13:WK} passes to the limits, and the Schwinger bounds and \Ezz\ of \cref{thm:s13:main}C(i) hold for $G^K_\infty$; moreover $G^K_\infty>0$ at the real points of $\Dtube$ and $G^K_J\le G^K_\infty$ at its cone points.
\end{theorem}

% s13_twk_proof: the proof of thm:s13:TWK (input from sec/s13.tex)
\begin{proof}
Throughout, $\cF=(f_1,\dots,f_n;h_1,\dots,h_p)$ is a family of real sources with supports in $Q_R$, $J\ge J_1(R)$, so that $MJ\ge R+\ell_c\ge R+10$: every source lies at $\ell^\infty$-distance $>1$ from $\R^3\setminus\CJ$, the support condition \eqref{eq:s04:S} holds for $\XJ$, and $J\ge J_0(\cF)$. For $\zeta\in\Dtube$ write $g=g(\zeta)$ and $k=k(\zeta)=k_+-k_-$ (\cref{def:s02:tube}); by \cref{lem:s02:E}(E3) the sampled sources satisfy $0\le2k_+\le s_1$ and $0\le2k_-\le H$ on $\Lk$. The inputs are: \Itwo\ with amplitude $s_1\in(0,\sigma_1)$ and constant $C_U$, from \Tthree\ and the torus inputs \cref{hyp:s03:T1,hyp:s03:T2} through \cref{prop:s10:U}, for $a\ne z$ with the constant $C_U+s_1\sup_k|z_k-a_k|$, finite by \cref{def:s03:adm} (\cref{hyp:s12:SCH}); \Ithree, which is \cref{prop:s10:RP}; \Ifour, which is \cref{lem:s02:BFS}; and the cell constant $B^{(a)}_{\rm cell}<\infty$ of the comparison convention (\cref{hyp:s12:Bcell}). Under these, \cref{cor:s11:W2prime} and \cref{prop:s11:W8} hold on the comparison convention and \cref{lem:s04:E}(d) on King's.

\emph{Step 1 (modulus).} \Cref{lem:s02:alg}(d) holds for every state of \cref{def:s02:acube}, killed sites included: $|Z^{a,K}_{k,J}(\zeta)|\le Z^{a,K}_{k,J}(\Re\zeta)$. It therefore suffices to bound $Z^{a,K}_{k,J}$ at the real point $\Re\zeta$, where $Z^{a,K}_{k,J}(\Re\zeta)=N^{a,K}(g,k)\defeq e^{-a_k\sS(k)}\avK{e^{F(g)+\Wick k}}$.

\emph{Step 2 (curvature off).} Let $\langle\cdot\rangle_{K,k}$ be King's cube with each single-site measure at an alive site multiplied by $e^{\eps^3k(x)r_x^2}$; it is a BFS measure (the factor keeps $\int e^{\alpha r^2}d\rho_x<\infty$ because $\lambda>0$, and the killed sites are unchanged), and $N^{a,K}(g,k)=N^{a,K}(0,k)\,\langle e^{F(g)}\rangle_{K,k}$. Its Gaussian is the Neumann photon Gaussian of $\XJ$, the same for both conventions, with covariance $C$ as in \cref{lem:s05:C1}. By \eqref{eq:s04:S} every plaquette with $g\ne0$ has its corners in $\XJ$, so $F(g)=\ip vA_k$ with $v=\dstar_kg$, and \cref{lem:s02:GD} (with $s=1$) and \cref{lem:s05:C1} give $\langle e^{F(g)}\rangle_{K,k}\le e^{\ip v{Cv}_k/2}\le e^{\norm g_{2,k}^2}$. This is the chain of \cref{thm:s05:CK}(a) in the presence of the real scalar source $k$. Hence $N^{a,K}(g,k)\le e^{\norm g_{2,k}^2}N^{a,K}(0,k)$.

\emph{Step 3 (the split).} \Cref{lem:s02:alg}(f) at $f=0$ with $k_1=k_+$, $k_2=k_-$: $N^{a,K}(0,k_+-k_-)\le N^{a,K}(0,2k_+)^{1/2}N^{a,K}(0,-2k_-)^{1/2}$.

\emph{Step 4 (the positive leg on King's convention).} Put $q\defeq2k_+$, a lattice function with $0\le q\le s_1$ supported in $Q_R\cap\Lk\subset\alive^K_J$. The observable $e^{\sum_x\eps^3q(x)r_x^2}$ is $\Theta_{0,q'}$ with $q'=\eps^3q\ge0$, of class $\cI$, and \eqref{eq:s04:S} holds for $\XJ$; Step~1 of the proof of \cref{lem:s04:G} (the kill step on $\cI$) gives $\avK{\Theta_{0,q'}}\le\ava{\Theta_{0,q'}}$, and the multiplier $e^{-(a_k+2\sW_\eps)\sS(q)}$ is one positive constant for both states (\cref{lem:s02:alg}(a)). Hence
\[
N^{a,K}_{k,\CJ}(0,q)\ \le\ N^a_{k,\CJ}(0,q)\ \le\ e^{\kappa\norm q_{1,k}}
\]
by \cref{cor:s11:W2prime}, whose hypotheses \Itwo, \Ithree\ hold (this is the cone instance of \cref{rem:s04:F}, as in \cref{rem:s11:Kconvention}). So $\frac12\ln N^{a,K}(0,2k_+)\le\kappa\norm{k_+}_{1,k}$.

\emph{Step 5 (the negative leg on King's convention).} Put $q\defeq2k_-$, a lattice function with $0\le q\le H$ supported in $Q_R$, and $J\ge J_1(R)$. \Cref{lem:s04:E}(d) gives $\ln N^{a,K}_{k,\CJ}(0,-q)\le K_-\wnorm q$ with the printed constant $K_-$ of \eqref{eq:s11:constants}, built from the cell constant $B^{(a)}_{\rm cell}$ of the comparison cell, which is finite by \cref{hyp:s12:Bcell}: by \cref{lem:s04:E}(c) every cell on which one of the cutoff sources $8\chi_sq$ does not vanish is a full cell inside $\CJ$ containing no killed site, so each factor of the wall insertion of \cref{lem:s04:E}(a) that carries a source is a comparison-cell functional, to which \cref{lem:s11:W4,lem:s11:W5} apply as printed, and the sum over cells is \cref{lem:s11:W67}; the scheme multiplier factorises over the cells (\cref{lem:s02:alg}(a)). The direction of the kill matters here and is the content of \cref{lem:s04:E}(e). Hence $\frac12\ln N^{a,K}(0,-2k_-)\le\frac12K_-\wnorm{2k_-}=K_-\wnorm{k_-}$.

\emph{Step 6 (assembly).} Steps 1--5 give, for every $k\ge0$, $J\ge J_1(R)$ and $\zeta\in\Dtube$,
\[
|Z^{a,K}_{k,J}(\zeta)|\le Z^{a,K}_{k,J}(\Re\zeta)\le\exp\big(\norm{g(\zeta)}_{2,k}^2+\kappa\norm{k_+(\zeta)}_{1,k}+K_-\wnorm{k_-(\zeta)}\big),
\]
which is \eqref{eq:s13:WK}. The constants are those of \eqref{eq:s11:constants}, with $B^{(a)}_{\rm cell}$ in $K_-$ and, for $a\ne z$, $C_U+s_1\sup_k|z_k-a_k|$ in $\kappa$.

\emph{Step 7 (the fixed-cube limits).} Each $Z^{a,K}_{k,J}$ is entire in $\zeta\in\C^{n+p}$: it is a finite-dimensional integral with Gaussian decay in $A$ ($\lmu>0$) and $e^{-\lambda\eps^3r^4}$ decay in $r$ at every alive site, the killed sites carrying no variables, and differentiation under the integral sign is justified. On a compact $Q\subset\Dtube$ the family $\{Z^{a,K}_{k,J}\}_{k\ge0}$ is uniformly bounded, by \eqref{eq:s13:WK} and \cref{lem:s02:E}(E1)--(E3): $\norm g_{2,k}^2\le C_0\fournorm g^2$, $\norm{k_+}_{1,k}\le C_0'\wnorm{k_+}$ and $\wnorm{k_\pm}\le\sum_j|\Re w_j|\wnorm{h_j}$, all continuous in $\zeta$. At every real point of $\Dtube$ the sequence converges: there $2\norm{k_+}_\infty<\sigma_+=s_1=s_W$ and the negative part is unrestricted, so \cref{hyp:s03:W1}, which is posed on King's convention (\cref{cor:k02:W1}(a)), gives the limit in $(0,\infty)$ for the reference scheme; for the admissible scheme $a$, $Z^{a,K}_{k,J}=e^{-(a_k-z_k)\sum_jw_j\sS_k(h_j)}Z^{z,K}_{k,J}$ with $a_k-z_k$ convergent and $\sS_k(h_j)\to\int h_j$ (\cref{def:s03:adm}, \cref{rem:k02:scheme}). The real points of $\Dtube$ form a non-empty open subset of $\R^{n+p}$ (\cref{def:s02:tube}), a set of uniqueness in the connected open $\Dtube$; by Montel's and Vitali's theorems (the fact (T3) stated before \cref{thm:s11:W}),$Z^{a,K}_{k,J}\to G^K_J$ locally uniformly on $\Dtube$, $G^K_J$ is holomorphic there, and \eqref{eq:s13:WK} passes to the limit with $\norm g_{2,k}\to\norm g_2$, $\norm{k_+}_{1,k}\to\norm{k_+}_{L^1}$ ($k_+$ is continuous with compact support) and the sampled $\wnorm{k_-}$ at most the continuum one (E3):
\begin{equation}\label{eq:s13:WKlim}
|G^K_J(\zeta)|\le\exp\big(\norm{g(\zeta)}_2^2+\kappa\norm{k_+(\zeta)}_{L^1}+K_-\wnorm{k_-(\zeta)}\big)\qquad(\zeta\in\Dtube,\ J\ge J_1(R)),
\end{equation}
a bound free of $J$. At real points $G^K_J>0$ (\cref{hyp:s03:W1}).

\emph{Step 8 (the cube limit on a sub-tube).} Split $h_j=h_j^+-h_j^-$ with $h_j^-\defeq\norm{h_j}_\infty\chi_j$, $\chi_j\in C^\infty_c$, $0\le\chi_j\le1$, $\chi_j=1$ on $\supp h_j$, $\supp\chi_j\subset Q_{R+1/2}$, and $h_j^+\defeq h_j+h_j^-\ge0$. For the split family $\cF''=(f;h^+,h^-)$ put $F^K_{k,J}(z,w^+,w^-)\defeq Z^{a,K}_{k,J}(\sum_iz_if_i,\sum_j(w_j^+h_j^++w_j^-h_j^-))$ and $\cD''\defeq\{(z,w^+,w^-)\in\C^{n+2p}:|\Re w_j^\pm|<\delta_j\}$ with $\sum_j\delta_j(\norm{h_j^+}_\infty+\norm{h_j^-}_\infty)<H/2$. On $\cD''$ the scalar source $k''$ of the split family has $\norm{k''_\pm}_\infty\le\norm{k''}_\infty<H/2\le\sigma_+/2$, so $\cD''$ lies in the tube of $\cF''$, whose supports lie in $Q_{R+1/2}$; Steps 1--7 for $\cF''$ and $J\ge J_1(R+\frac12)$ give $F^K_{k,J}\to F^K_J$ locally uniformly on $\cD''$ with the $J$-free bound \eqref{eq:s13:WKlim} for $\cF''$. On $\Omega''_+\defeq\cD''\cap((\mathrm i\R)^n\times(0,\infty)^{2p})$ the curvature coefficients are imaginary and the total scalar source $\sum_j(w_j^+h_j^++w_j^-h_j^-)$ is nonnegative, so the observable is of class $\cI$ and \eqref{eq:s04:S} holds for $\XJ$ ($MJ\ge R+\frac12+\ell_c$); \cref{cor:s04:G1} gives $F^K_{k,J}\le F^K_{k,J+1}$ at every $k$ (the scheme multiplier is common, \cref{lem:s02:alg}(a)), and $k\to\infty$ gives $F^K_J\le F^K_{J+1}$ on $\Omega''_+$. Being nondecreasing and bounded there, $J\mapsto F^K_J$ converges on $\Omega''_+$, a non-empty open subset of the maximal totally real subspace $(\mathrm i\R)^n\times\R^{2p}$; $\cD''$ is convex, hence connected, and the Montel--Vitali fact (T3) stated before \cref{thm:s11:W} gives $F^K_J\to F^K_\infty$ locally uniformly on $\cD''$. At $w^+_j=w_j$, $w^-_j=-w_j$ the total scalar source is $\sum_jw_j(h_j^+-h_j^-)=\sum_jw_jh_j$, so $F^K_{k,J}(z,w,-w)=Z^{a,K}_{k,J}(\zeta)$ exactly at every $k$; hence $G^K_J\to G^K_\infty$ locally uniformly on $\cD'\defeq\{\zeta:|\Re w_j|<\delta_j\}$, which lies in $\Dtube$ because there $\norm{k_\pm}_\infty\le\norm k_\infty<\sum_j\delta_j\norm{h_j}_\infty<H/2$.

\emph{Step 9 (the whole tube).} The family $\{G^K_J\}_{J\ge J_1(R+1/2)}$ is locally uniformly bounded on $\Dtube$ by \eqref{eq:s13:WKlim}; by Montel's theorem every subsequence has a further subsequence converging locally uniformly to a holomorphic function, any two such limits agree on the open set $\cD'$ by Step~8, and $\Dtube$ is connected, so all subsequential limits coincide: $G^K_J\to G^K_\infty$ locally uniformly on $\Dtube$, $G^K_\infty$ is holomorphic, and \eqref{eq:s13:WKlim} passes to $G^K_\infty$. The finitely many indices $J_1(R)\le J<J_1(R+\frac12)$ do not affect the limit.

\emph{Step 10 (positivity and the cone).} For a real $\zeta\in\Dtube$ pick $\epsilon>0$ with $-\epsilon\zeta\in\Dtube$. The function $\varphi_J(t)\defeq\ln G^K_J(t\zeta)$ is finite and convex on $[-\epsilon,1]$ (a pointwise limit of the convex functions $\ln Z^{a,K}_{k,J}(t\zeta)$, \cref{lem:s02:alg}(d), with positive limits at real points) and $\varphi_J(0)=0$, so $\varphi_J(1)\ge-\varphi_J(-\epsilon)/\epsilon\ge-C_\epsilon/\epsilon$, $C_\epsilon$ the exponent of \eqref{eq:s13:WKlim} at $-\epsilon\zeta$; hence $G^K_\infty(\zeta)\ge e^{-C_\epsilon/\epsilon}>0$. At the cone points of $\Dtube$ ($z\in(\mathrm i\R)^n$, $w\in\R^p$, $k(\zeta)\ge0$), \cref{cor:s04:G1} at every $k$ and $k\to\infty$ give $G^K_J\le G^K_{J+1}\le G^K_\infty$.

\emph{Step 11 (the Schwinger bounds and \Ezz).} Put $S^K_{n,p}(f;h)\defeq\partial^{(n)}_z\partial^{(p)}_wG^K_\infty(0)$, $\sigma_0=H/4$, $K'=\kappa\pi^2+K_-$. Normalise $\norm{f_i}_2=1$ and let $\rho_j\defeq\sigma_0/(p\wnorm{h_j})$. On the polydisc $\{|w_j|\le\rho_j\}\times\C^n$ we have $\norm k_\infty\le\sum_j|\Re w_j|\norm{h_j}_\infty\le\sum_j\rho_j\wnorm{h_j}=\sigma_0<H/2\le\sigma_+/2$ (\cref{lem:s02:E}(E2)), so the polydisc lies in $\Dtube$, where $G^K_\infty$ exists with the bound \eqref{eq:s13:WKlim}, whose exponent is at most $\norm{g(z)}_2^2+\kappa\pi^2\wnorm{k_+}+K_-\wnorm{k_-}\le\norm{g(z)}_2^2+K'\wnorm k\le\norm{g(z)}_2^2+K'\sigma_0$ (E2). Cauchy's estimate for the mixed first derivative in the $w$'s at $w=0$, $z$ fixed, gives $|\partial^{(p)}_wG^K_\infty(z,0)|\le\prod_j\rho_j^{-1}\sup|G^K_\infty|\le(p/\sigma_0)^p\prod_j\wnorm{h_j}\,e^{K'\sigma_0}e^{\norm{g(z)}_2^2}$, and $p^p\le e^pp!$ (E4). The function $z\mapsto\partial^{(p)}_wG^K_\infty(z,0)$ is entire ($\C^n\times\{w\}\subset\Dtube$ for $|\Re w_j|$ small) and bounded by $A\,e^{\norm{g(z)}_2^2}$ with $A\defeq e^{K'\sigma_0}(e/\sigma_0)^pp!\prod_j\wnorm{h_j}$; on the polydisc of radius $r=(2n)^{-1/2}$ about $0$, $\norm{g(z)}_2\le\sum_i|\Re z_i|\le nr$, so Cauchy's estimate for the mixed first $z$-derivative gives $|S^K_{n,p}|\le A\,r^{-n}e^{n^2r^2}=A\,(2n)^{n/2}e^{n/2}\le A\,(\sqrt2e)^n(n!)^{1/2}$ by (E4). Restoring the normalisation by multilinearity,
\[
|S^K_{n,p}(f_1,\dots,f_n;h_1,\dots,h_p)|\le e^{K'\sigma_0}(\sqrt2e)^n(n!)^{1/2}(e/\sigma_0)^p\,p!\prod_i\norm{f_i}_2\prod_j\wnorm{h_j}.
\]
By (E2), $\norm{f_i}_2\le\frac\pi{2\sqrt2}\fournorm{f_i}$ and $\wnorm{h_j}\le\fournorm{h_j}$, and $(n!)^{1/2}p!\le(n+p)!$, so the right side is at most $\sigma_{n+p}\prod_i\fournorm{f_i}\prod_j\fournorm{h_j}$ with $\sigma_N=e^{K'\sigma_0}C_1^NN!$, $C_1=2\max(e\pi/2,e/\sigma_0)$; for complex test functions each argument is split into real and imaginary parts, which gives the factor $2^{n+p}\le2^N$ absorbed in $C_1$. The $S^K_{n,p}$ are therefore continuous multilinear forms in the norms $\fournorm\cdot$, extend to $\cS(\R^3;\Lambda^2_\C)^n\times\cS(\R^3;\C)^p$ and are tempered distributions, as in \cref{thm:s05:C}(c), and $C_1^N/(N!)^{\beta-1}\to0$ for every $\beta>1$: this is \Ezz\ with $s=4$ and every $\beta>1$, with the constants of \cref{thm:s13:main}C(i).
\end{proof}


\begin{proposition}[the Wick sector on King's convention]\label{prop:s13:WK}
Let the couplings satisfy \Tthree\ and let $a$ be an admissible scheme. Then, with the constants of \cref{thm:s13:main}C:
\begin{enumerate}[label=\textup{(\alph*)},leftmargin=*,nosep]
\item $G^K_J\defeq\lim_kZ^{a,K}_{k,J}$ exists locally uniformly on $\Dtube$ for every $J\ge J_1(R)$, $G^K_\infty\defeq\lim_JG^K_J$ exists locally uniformly on $\Dtube$, both are holomorphic, $|G^K_\infty(\zeta)|\le\exp(\norm{g(\zeta)}_2^2+\kappa\norm{k_+(\zeta)}_{L^1}+K_-\wnorm{k_-(\zeta)})$, and the Schwinger bounds and \Ezz\ of \cref{thm:s13:main}C(i) follow.
\item $G^K_\infty(\tau_bR\cdot\zeta)=G^K_\infty(\zeta)$ on $\Dtube$ for all $b\in\R^3$, $R\in\Oh$, $W^a$ a scalar.
\item For every family $\cF'$ with nonnegative scalar members the Wick-tube object $Z^{a,K}_{\infty,T}\defeq\lim_J\lim_kZ^{a,K}_{k,J}$ exists locally uniformly on $T_{\cF'}$ and equals $G^K_\infty$ on $T_{\cF'}\cap\Dtube$; reflection positivity about every coordinate plane holds in the matrix form at amplitude $<\sigma_W$ and in the polynomial form, and \OSth, for $Z^{a,K}_{\infty,T}$, $G^K_\infty$ and their Schwinger functions.
\item The measure and weak \OSf\ of \cref{thm:s13:main}C(v),(vi) for the joint field under $\NDK$, with the local convergence of the laws of $\chi(F_{\Lk},W^a_{\Lk})$ under $\NDK_{k,J}$ along the full sequence $k\to\infty$ at every fixed $J\ge\max(J_0(\chi),J_1(R_\chi))$, then $J\to\infty$.
\end{enumerate}
\end{proposition}

\begin{proof}
The proofs of \cref{thm:s11:W}(i), \cref{lem:s11:TWtube}, \cref{thm:s06:S1K} with scalar sources, \cref{thm:s08:OS2WT} with \cref{prop:s08:T}, and \cref{thm:s09:S2K} in the Wick sector are run on $\NDK$ with the replacements below; every other step of theirs is convention-free, since it sees the states only as members of the class of \cref{lem:s02:BFS} or concerns the limit object alone.

(a) is \cref{thm:s13:TWK}, whose proof contains the fixed-cube limits (Step~7), the cube limit on the connected tube with the cube monotonicity of \cref{cor:s04:G1} in place of the first inequality of \cref{lem:s02:cut} (Steps 8--9), the positivity at real points and the cone monotonicity (Step~10), and the Schwinger bounds by Cauchy's estimates on a polydisc in the tube (Step~11).

(b) The steps of \cref{thm:s06:S1K} with scalar sources. Coordinate permutations are exact on $\NDK$ for $(f,h)$ by \cref{lem:s06:PK}. The reflection sandwich is \cref{lem:s04:C}(b) at the cone points of the tube: for $(f,g)\in\cK$ with $\norm g_\infty<\sigma_+/2$ the point $(\mathrm i,1)$ lies in $\Dtube$, so (a) supplies the three $k$-limits and the two $J$-limits of the two-step sandwich $C^K_{J\mp2}$. Conjugation and group closure on the cone are unchanged. Off the cone, the identity theorem on the split family's tube of \cref{lem:s11:TWtube}, where (a) gives the limits, extends the identity from the cone points to the connected tube; so \cref{lem:s04:C} is needed on $\cK$ only, and no class-$\cD$ form of it enters (the proof splits $h=h^+-h^-$ and never applies the sandwich to a negative scalar source). Translations are \cref{lem:s06:TK}, stated on $\cK$ and applied at the cone points of the tube with (a) for the limits, then the continuity route of \cref{cor:s06:TK} with the Wick-sector second-moment bound, which Cauchy's estimates give from \eqref{eq:s11:W} on $\NDK$; the translated family needs $J\ge\max(J_0+1,J_1(R+M+\frac12))$. Invariance about every point and differentiation concern the limit object.

(c) \Cref{thm:s08:OS2WT} and \cref{prop:s08:T} on $\NDK$: in the sandwich of the $J$-limit the upper box is \cref{lem:s04:D} (the printed boxes $\CJ$, $C_{J+2n}$) and the lower end carries the kill step of \cref{rem:s04:F}; the positive-leg bound \cref{cor:s11:W2prime} holds on $\NDK$ by \cref{rem:s04:F} prepended to \cref{prop:s11:W1prime} (\cref{rem:s11:Kconvention}); the cube monotonicity in both classes is \cref{cor:s04:G1}; the fixed-cube limits on $T_{\cF'}$ are \cref{thm:k02:W1}(i), whose domain $\{\sup_x2k_1<s_W\}$, $s_W=s_1$, contains $T_{\cF'}$ (for the reference scheme, hence for $a$ by \cref{def:s03:adm}), and on $\Dtube$ they are (a); the cube limit on $T_{\cF'}$ is then Montel's theorem with the uniqueness set of \cref{def:s08:enlarged}(iii), as in \cref{prop:s08:T}; the translations to the planes $c\notin D$ are (b). \OSth\ and the polynomial form use only the existence and holomorphy of the limit near $0$ and charge conjugation (\cref{prop:s08:OS3,prop:s08:poly}).

(d) The steps of \cref{thm:s09:S2K} on $\cV=\Lambda^2\oplus\R$ with the scheme-shifted field: \eqref{eq:s11:W} on $\NDK$ in the moment bounds, (a) for the limit functionals and the lattice laws, (b) and (c) for the invariance and the reflection positivity consumed by the abstract chain of \cite{SymPhase}. The second moments of the lattice field under $\NDK_{k,J}$ come from \eqref{eq:s11:W} on $\NDK$ by Cauchy's estimates, so tightness and the identification of the limit law run as in the curvature sector, with the fixed-cube limits of (a) along the full sequence, for $J\ge\max(J_0(\chi),J_1(R_\chi))$; the remaining steps concern the limit measure.
\end{proof}

% s13_limk: the Wick-sector LIM states of the comparison family are King's tube object (input from sec/s13.tex)
\begin{proposition}[the comparison family on the Wick tube]\label{prop:s13:LIMK}
Let the couplings satisfy \Tthree\ and let $a$ be an admissible scheme.
\begin{enumerate}[label=\textup{(\alph*)},leftmargin=*,nosep]
\item Every increasing sequence of spacings has a continuum subsequence $\cK$ for the Wick sector of the comparison family (\cref{def:s12:contsub}).
\item For every such $\cK$, every family $\cF$ of real sources with supports in $Q_R$ and every $\zeta\in\Dtube$, the LIM state of \cref{def:s12:LIMstate} is King's tube object:
\[
Z^{a,\cK}_\infty(\zeta)\ \defeq\ \lim_{J\to\infty}\lim_{m\to\infty}Z^a_{k_m,J}(\zeta)\ =\ G^K_\infty(\zeta)\qquad(\zeta\in\Dtube).
\]
In particular the Wick-sector LIM state does not depend on $\cK$, and the cube limit of every subsequential fixed-cube limit of the comparison family on the tube is $G^K_\infty$.
\end{enumerate}
\end{proposition}

\begin{proof}
(a) By \cref{hyp:s12:Bcell,hyp:s12:SCH} the scheme has the cell bound and the scheme bound, so \cref{lem:s12:LIM0}(b) applies.

(b) Fix $\cK=(k_m)$ and $\cF$. By \cref{lem:s12:LIM1} and the Wick sector of \cref{lem:s12:LIM0}(b), $Z^{a,\cK}_\infty$ exists locally uniformly on $\Dtube$ and is holomorphic there; by \cref{thm:s13:TWK} so is $G^K_\infty$. Take the split family $\cF''=(f;h^+,h^-)$ of Step~8 of the proof of \cref{thm:s13:TWK}, with supports in $Q_{R+1/2}$, the domain $\cD''$ and the set $\Omega''_+=\cD''\cap((\mathrm i\R)^n\times(0,\infty)^{2p})$ defined there; write $F^K_{k,J}$ and $F^{(a)}_{k,J}$ for the functionals of $\cF''$ on King's and on the comparison cube. Let $J$ be so large that $MJ>R+\frac12+\ell_c$. At a point of $\Omega''_+$ the curvature coefficients are purely imaginary and the total scalar source is nonnegative, so the observable is of class $\cI$, and the support condition \eqref{eq:s04:S} holds for $\XJ$; \cref{lem:s04:G} (with its clause for the scheme functionals) gives, at every $k$,
\[
F^K_{k,J}\ \le\ F^{(a)}_{k,J}\ \le\ F^K_{k,J+1}\qquad\text{on }\Omega''_+ .
\]
Along $\cK$ the two ends converge, by \cref{thm:s13:TWK} for the family $\cF''$ (whose fixed-cube limits exist along the full sequence, hence along $\cK$), to the fixed-cube limits $F^K_J$, $F^K_{J+1}$ of Step~8 of that proof, and the middle term converges to the fixed-cube limit $F^{(a),\cK}_J$ of the comparison family along $\cK$ (\cref{def:s12:contsub} for $\cF''$, whose tube contains $\cD''$ and whose threshold is $MJ>R(\cF'')+\ell_c$, $R(\cF'')\le R+\frac12$). Letting $J\to\infty$, the two ends tend to $F^K_\infty$ (Steps 8--9 of the proof of \cref{thm:s13:TWK}) and the middle term to the LIM state $F^{(a),\cK}_\infty$ of $\cF''$ (\cref{lem:s12:LIM1}). Hence $F^{(a),\cK}_\infty=F^K_\infty$ on $\Omega''_+$, a non-empty open subset of the maximal totally real subspace $(\mathrm i\R)^n\times\R^{2p}$, and both functions are holomorphic on the connected open set $\cD''$; by the identity theorem (the fact (T2) stated before \cref{thm:s11:W}) they agree on $\cD''$. At $w^+_j=w_j$, $w^-_j=-w_j$ the lattice functionals of $\cF''$ coincide exactly with those of $\cF$ at $\zeta$, on both conventions and at every $k$ (equal total sources); so $Z^{a,\cK}_\infty=G^K_\infty$ on $\cD'=\{\zeta:|\Re w_j|<\delta_j\}\subset\Dtube$. Both are holomorphic on the connected $\Dtube$ and agree on the open set $\cD'$, hence on $\Dtube$. The last clause is \cref{def:s12:LIMstate}: every subsequential fixed-cube limit on the tube is the fixed-cube limit along some continuum subsequence.
\end{proof}


\begin{remark}[the three amplitude constants]\label{rem:s13:amplitudes}
Three amplitudes occur in the Wick sector, all built from $s_1$: the matrix form of \OSt\ is proved at amplitude $\amp(\zeta_i)<\sigma_W=\sigma_+/8$ on the Wick tube (\cref{thm:s08:OS2WT}); the tube $\Dtube$ has half-widths $\sigma_+/2$ for the positive and $H/2=\sigma_+/32$ for the negative part of the real Wick source, and the form of \OSt\ proved for LIM states (\cref{thm:s12:LIM}) takes $\amp(\zeta_i)<H/2<\sigma_W$, so that every Gram point lies in $\Dtube$, where those limits exist; and $\sigma_0=H/4$ is the radius of the polydisc in Cauchy's estimates for \Ezz. The consumer of \OSt\ in clauses (v),(vi), the reflection positivity of exponentials $e^{\mathrm i\Phi(u)}$ with real $u$ in the clustering theorems of \cite{SymPhase}, needs purely imaginary Wick coefficients only, that is amplitude $0$ (\cref{rem:s08:amplitudes}).
\end{remark}

\begin{remark}[the photon mass and the couplings]\label{rem:s13:mu0}
The photon mass $\lmu>0$ is fixed throughout \cref{thm:s13:main}. No relation $\lmu\le m$ is used or needed: the one input whose printed proof in the torus paper \cite{SymPhase} carries the restriction $\lmu\in(0,\mus]$, the lattice reflection positivity of King's tori, holds for all $\lmu>0$, since that proof uses no bound on $\lmu$ and the paper's own lemma on the finite-volume photon covariance lifts the scope (\cref{prop:s10:RP}, \cref{lem:s14:CW}(d)). Reflection positivity of the cube limit needs the margin $R(\lmu)=\pi\sqrt3/(2\lmu)$ of \eqref{eq:s07:margin} inside the upper half of the sandwich (\cref{lem:s07:U}) and nowhere in the conclusion. The constants of Leg~C are absolute; those of Leg~W depend on $\lmu$ and are not uniform in it. The couplings range over all of $e>0$, $\lambda,m,\lmu>0$: their quantifier is the sources' (\cref{lem:s14:CW}(e), \cref{rem:s14:quantifier}), the stopping scale of the lattice constructions behind \cref{hyp:s03:EC,hyp:s03:W1,hyp:s03:T1} being chosen after the point and the datum, with the mass cap of \cite[Prop.~3.1]{Ba2} one more one-sided cap on that scale (\cref{rem:s14:BALm}); so the theorem contains no sentence ``$e,\lambda$ sufficiently small'', and the only coupling condition of the Wick sector is the one-sided sign clause \Tthree\ (\cref{lem:s14:CW}(b)), proved on the explicit regions of \cref{prop:s14:T3}.
\end{remark}
